\documentclass[12pt,a4paper]{article}

% Packages
\usepackage[utf8]{inputenc}
\usepackage[T1]{fontenc}
\usepackage{amsmath,amssymb}
\usepackage{graphicx}
\usepackage{booktabs}
\usepackage{float}
\usepackage{geometry}
\usepackage{setspace}
\usepackage{natbib}
\usepackage{hyperref}

\geometry{margin=2.5cm}
\onehalfspacing

\title{
Dynamic Volatility and Market Risk Forecasting:
A Comparative Analysis of Econometric and Machine Learning Models
with a Hybrid Extension
}

\author{Alejandro Basabe Oleaga}

\date{2026}

\begin{document}

\maketitle

\begin{abstract}
This thesis compares econometric (EWMA, GARCH, GJR-GARCH), machine learning (Random Forest,
XGBoost), and hybrid volatility-forecasting models for the S\&P 500 (2011--2026), asking not only
which model forecasts most accurately but whether that accuracy translates into better
market-risk measurement.

Using a strictly out-of-sample walk-forward framework with Value-at-Risk backtesting, machine
learning is found to offer no statistically significant forecasting advantage over the
asymmetric GJR-GARCH model, and performs significantly worse under a loss function that
specifically penalises underprediction of risk. A hybrid model combining GJR-GARCH with a
machine-learning residual correction achieves the most accurate point forecasts of any model
considered, a statistically confirmed result -- yet it fails Value-at-Risk backtesting at the
99\% confidence level exactly as badly as the machine learning models its correction is drawn
from. This disconnect traces to a specific mechanism: the correction learns a near-uniform
downward adjustment that improves average accuracy while shading forecasts down hardest on
precisely the days that matter most for risk. A subsequent attempt to correct this directly, by
training the residual model with a loss function aligned with the definition of Value at Risk
itself, fails for a different and more general reason: a correction calibrated on a historically
volatile period does not transfer to a calmer one.

The central finding of this thesis is that improvements in standard forecast-accuracy metrics do
not necessarily translate into improvements in market-risk measurement, and can coincide with a
statistically confirmed deterioration in risk calibration. GJR-GARCH alone remains the
best-calibrated model for market-risk purposes throughout this thesis, and no combination or
correction attempted here improves on it in that respect.
\end{abstract}

\noindent\textbf{Keywords:} GARCH; GJR-GARCH; volatility forecasting; machine learning;
Value at Risk; hybrid models

\tableofcontents
\newpage

\section{Introduction}
% ============================================================
\subsection{Motivation}
% ============================================================
Financial market volatility is an important component of quantitative finance and risk management.
Accurate volatility forecasts are relevant for the measurement of market risk, portfolio management
and financial decision-making. Traditional econometric models, particularly models from the ARCH
and GARCH families, have been widely used to model the time-varying nature of financial volatility.

In recent years, machine learning methods have become increasingly relevant in financial forecasting.
Their ability to model nonlinear relationships provides a potential alternative to traditional
econometric approaches. However, greater predictive flexibility does not necessarily imply better
out-of-sample forecasts or more accurate measures of financial risk.

This work therefore compares traditional econometric volatility models with machine learning models
using a common out-of-sample forecasting framework. The analysis is subsequently extended to a
hybrid model combining econometric and machine learning approaches. The final objective is to
determine whether improvements in volatility forecasting translate into improvements in market-risk
measurement.
% ============================================================
\subsection{Research Question}
% ============================================================
This thesis addresses the two following questions:

\begin{enumerate}
    \item Can machine-learning methods improve the out-of-sample forecasting of financial market volatility compared with traditional econometric models, and does any improvement translate into more accurate market-risk measures?
    \item Does combining GARCH and machine learning provide additional forecasting value relative to the individual models?
\end{enumerate}
% ============================================================
\subsection{Objectives}
% ============================================================
The main objective of this work is to evaluate whether machine-learning models provide improvements
in the out-of-sample forecasting of financial market volatility relative to traditional econometric
models, and whether these differences translate into more accurate market-risk measures.

A secondary objective is to develop a hybrid GARCH--machine learning model and determine whether
combining the two approaches provides additional forecasting value compared with the corresponding
individual models.

The objective is not to assume that machine learning models are superior, but to evaluate their
performance under a common and strictly out-of-sample framework.
% ============================================================
\subsection{Contributions}
% ============================================================
The main contribution of this work is an empirical comparison between traditional econometric
volatility models and machine learning models using the same financial data, forecasting horizon
and out-of-sample evaluation criteria.

The analysis is complemented by a market-risk application through the estimation and backtesting
of Value at Risk (VaR) and Expected Shortfall (ES). Finally, a hybrid GARCH--machine learning model
is developed to evaluate whether the combination of econometric and machine learning information
provides an additional improvement in forecasting performance.
% ============================================================
\subsection{Structure of the Thesis}
% ============================================================
Following this introduction, the literature review presents the main theoretical and empirical
references related to volatility modelling, machine learning forecasting, hybrid models and
financial risk measurement.\\
The data section describes the financial dataset, the construction of returns and the definition of the volatility forecasting target.\\
The methodology section describes the econometric models, machine learning models, hybrid model, forecast evaluation criteria and market-risk backtesting procedures.\\
The results section presents and compares the out-of-sample forecasting performance of the different models and analyses their implications for VaR and Expected Shortfall.\\
Finally, the conclusions summarise the main findings and answer the research questions proposed at the beginning of the work.
% ============================================================
\section{Literature Review}
% ============================================================
The literature relevant to this thesis can be organised into four themes: the
econometric foundations of volatility modelling, the measurement of market risk,
the evaluation of volatility forecasts, and the more recent application of machine
learning and hybrid approaches to the same problem.

The econometric treatment of time-varying volatility originates with
\citet{engle1982}, who introduced the ARCH model to capture the clustering of
volatility observed in macroeconomic and financial time series.
\citet{bollerslev1986} generalised this framework with the GARCH model, allowing
conditional variance to depend on its own lagged values in addition to lagged
squared shocks. Together these two papers established conditional
heteroskedasticity models as the standard econometric benchmark for volatility
forecasting, and GARCH(1,1) in particular remains difficult to outperform
consistently in out-of-sample comparisons. This thesis extends the baseline GARCH
specification with an asymmetric GJR-GARCH variant, motivated by the
well-documented tendency of volatility to react more strongly to negative than to
positive returns.

Once volatility is forecast, it must be translated into a usable measure of
market risk. \citet{jorion1997} formalised Value at Risk as an industry and
regulatory benchmark, while \citet{artzner1999} showed that VaR lacks
subadditivity and is therefore not a coherent risk measure, motivating the use of
Expected Shortfall as a complementary, coherent alternative. This thesis follows
that recommendation by reporting both VaR and Expected Shortfall rather than VaR
alone.

A separate strand of the literature addresses how volatility forecasts themselves
should be evaluated. \citet{poon2003} review the difficulties inherent in this
exercise, in particular the sensitivity of forecast rankings to the choice of
realised-volatility proxy and evaluation loss function, and stress the importance
of genuinely out-of-sample testing. This motivates the use of multiple loss
functions (MAE, RMSE and QLIKE) in the present work rather than a single metric,
since no individual measure is uniformly preferred.

More recently, machine learning methods have been applied to volatility
forecasting as an alternative to purely econometric specifications.
\citet{gunnarsson2024} survey this literature and report that machine learning
methods for volatility forecasting are generally effective, though traditional
econometric models continue to perform comparably in many cases; they identify
hybrid specifications, combining machine learning with econometric structure, as
a particularly promising direction for future research --- the approach adopted
later in this thesis's own hybrid extension. \citet{misra2025} provide a more
direct comparison between GARCH, XGBoost and LSTM models for S\&P 500 volatility
forecasting over 2018--2024 and report that both machine learning models
outperform GARCH by a substantial margin on $R^2$, RMSE and MAPE. That
comparison is, however, a working paper rather than a peer-reviewed study, and
it evaluates forecasts only against symmetric accuracy criteria. Since
\citet{poon2003} show that forecast rankings are themselves sensitive to the
loss function chosen, results established on symmetric criteria alone do not
settle whether machine learning also improves the risk-relevant properties of a
forecast. That gap, and the absence of a uniform conclusion in favour of machine
learning across this broader literature, motivates the neutral, criteria-driven
comparison adopted in this thesis rather than an a priori assumption of
superiority in either direction.

A smaller body of work combines the two approaches directly.
\citet{kartsonakis2021} compare standalone GARCH and neural network models
against a hybrid specification for the Italian stock market and report gains
from the hybrid approach, though without a fully out-of-sample evaluation
matched across models. This result motivates the hybrid extension developed
later in this thesis, in which a machine learning model is trained specifically
on the residual information left unexplained by GARCH, rather than combined
with it in an unstructured way.

Taken together, this literature establishes GARCH-family models as the relevant
econometric benchmark, VaR and Expected Shortfall as complementary risk measures,
and highlights both the promise and the methodological pitfalls of applying
machine learning to volatility forecasting. What remains comparatively unexplored
is a single, strictly out-of-sample comparison that carries econometric,
machine-learning and hybrid volatility forecasts through to their implications
for backtested market-risk measures --- which is the gap this thesis addresses.
This finding carries a practical implication beyond the specific models compared here.
Selecting a volatility model on the basis of average point-forecast accuracy alone, without
separately evaluating the risk measures built on top of it, risks adopting a model that appears
superior on the criterion used to choose it while performing worse on the criterion that
actually matters for risk management. The two should be assessed independently, using metrics
matched to their eventual use, rather than assumed to move together.
% ============================================================
\section{Data}
% ============================================================
\subsection{Dataset}
% ============================================================
The analysis uses approximately fifteen years of daily S\&P 500 data, covering the period from
January 2011 to 7 July 2026. Data is extracted using the \texttt{yfinance} package (ticker
\texttt{\^{}GSPC}) and the end date is fixed at 7 July 2026 so that the dataset (and therefore all
downstream results) remains identical across every re-run of the pipeline.

The dataset contains daily Open, High, Low, Close and Volume observations. Adjusted closing prices
are used for the calculation of daily returns; since the S\&P 500 is a price index rather than a
total-return index, adjusted and unadjusted closing prices coincide, and no dividend adjustment is
required at the index level.

The raw data is cleaned by removing duplicate entries and any weekend dates accidentally included
by the data source. Rather than deleting anomalous records -- such as days with missing information,
zero trading volume, or extreme market swings -- the cleaning script preserves these rows and flags
them for review, so that no information is silently discarded. Daily logarithmic returns are then
calculated to prepare the data for modelling.

Two flagged events were identified in the sample:

\begin{itemize}
    \item \textbf{16 March 2020} (log return $-12.8\%$): corresponds to the COVID-19 market crash
    and reflects a genuine market event rather than a data error.
    \item \textbf{24 May 2023} (zero trading volume): likely a data-provider reporting quirk; the
    price data for this day appears otherwise consistent.
\end{itemize}

After cleaning, the dataset comprises 3,899 daily observations.
% ============================================================
\subsection{Returns}
% ============================================================
Daily logarithmic returns are calculated as

\begin{equation}
r_t = \ln\left(\frac{P_t}{P_{t-1}}\right),
\end{equation}
where $P_t$ represents the adjusted closing price of the S\&P 500 on day $t$. Figure~\ref{fig:returns_timeseries}
plots the resulting series over the full sample.
\begin{figure}[H]
\centering
\includegraphics[width=0.85\textwidth]{returns_timeseries.pdf}
\caption{Daily logarithmic returns of the S\&P 500, 2011--2026. The 16 March 2020 COVID-19 crash is
annotated; periods of high and low volatility cluster together rather than occurring randomly,
motivating the use of conditional heteroskedasticity models.}
\label{fig:returns_timeseries}
\end{figure}

% ============================================================
\subsection{Volatility Target}\label{sec:target}
% ============================================================
The dependent variable for the forecasting exercise is a five-day close-to-close realised
volatility measure.

At each date $t$, the realised variance over the following five trading days is calculated as

\begin{equation}
RV_{t+1:t+5} =
\sum_{i=1}^{5} r_{t+i}^{2}.
\end{equation}

The corresponding realised volatility is defined as

\begin{equation}
\sqrt{\frac{252}{5}RV_{t+1:t+5}},
\end{equation}

where the factor $252$ annualises the volatility.\\
This measure provides a less noisy forecasting target than the squared return of a single trading
day while allowing the analysis to be performed entirely using daily data.
% ============================================================
\section{Methodology}
% ============================================================
The methodology is designed to provide a common and strictly out-of-sample framework for comparing
econometric and machine learning approaches.

All models forecast the same five-day volatility target using only information available at
the forecasting date.

\subsection{Out-of-Sample Evaluation}\label{sec:oos}

The dataset is divided chronologically into a training period and an out-of-sample testing period:
observations from 3 January 2011 to 30 December 2022 (3,020 observations) are used for model
development, and observations from 3 January 2023 to 7 July 2026 (879 observations) are reserved
for the final out-of-sample evaluation. Because the five-day-ahead target requires five subsequent
trading days of returns, the final five observations of the training period lack a valid label;
these are excluded from the estimation of any model trained directly on the target variable (3,015
usable training examples for such models), while models that require only return data, such as the
econometric models, are unaffected.

A walk-forward expanding-window procedure is used rather than a random train-test split. At
each forecasting date, the model is estimated using information available up to that date and
subsequently produces a forecast for the following five trading days.

This procedure avoids look-ahead bias and provides a realistic assessment of the models' forecasting
performance.

\subsection{Econometric Models}

Three econometric approaches are considered.

\subsubsection{EWMA}

Exponentially Weighted Moving Average volatility is used as a simple benchmark, with decay
parameter $\lambda = 0.94$, the standard convention introduced by the RiskMetrics methodology. EWMA
gives greater weight to recent observations and provides a straightforward representation of
volatility clustering.

\subsubsection{GARCH(1,1)}

The standard GARCH(1,1) model is used as the main symmetric econometric benchmark. Its
conditional variance is specified as

\begin{equation}
\sigma_t^2 =
\omega +
\alpha \epsilon_{t-1}^2 +
\beta \sigma_{t-1}^2,
\end{equation}

where $\omega$, $\alpha$ and $\beta$ are model parameters.

\subsubsection{GJR-GARCH}

A GJR-GARCH model is included to account for potential asymmetry in the response of volatility
to positive and negative returns. This allows the analysis to determine whether explicitly modelling
asymmetric volatility improves forecasting performance relative to the standard GARCH model.

For comparability, all three econometric models are converted from a one-day-ahead conditional
variance forecast into the same five-day annualised volatility scale as the forecasting target
(Section~\ref{sec:target}), using an identical annualisation procedure applied uniformly across
models. GARCH(1,1) and GJR-GARCH produce a full five-step-ahead conditional variance path, which is
allowed to mean-revert towards its long-run level over the forecast horizon. EWMA, which has no
mean-reversion mechanism, instead applies its single-step variance forecast uniformly across all
five days -- the standard treatment for EWMA over a multi-day horizon. The three models therefore
differ not only in their conditional-variance dynamics but also in how that dynamic is extended
across the forecast horizon, a distinction worth bearing in mind when interpreting differences in
forecasting accuracy.

\subsection{Machine Learning Models}\label{sec:ml_features}

Two machine learning algorithms are considered: Random Forest and XGBoost.

These models are selected because they can capture nonlinear relationships while remaining
substantially simpler to implement and interpret than more complex deep-learning architectures.

To keep the comparison against the econometric models uncontaminated, the machine learning models
used to answer Research Question 1 are trained using a limited set of six features constructed
exclusively from returns, without reference to any econometric model's own output:

\begin{table}[H]
\centering
\begin{tabular}{ll}
\toprule
\textbf{Feature} & \textbf{Definition} \\
\midrule
Same-day return & $r_t$ \\
Absolute return & $|r_t|$ \\
Squared return & $r_t^2$ \\
5-day historical volatility & $\sigma_{5,t}$ \\
10-day historical volatility & $\sigma_{10,t}$ \\
20-day historical volatility & $\sigma_{20,t}$ \\
\bottomrule
\end{tabular}
\caption{Features used by the machine learning models for the Research Question 1 comparison.}
\end{table}

Table~\ref{tab:ml_hyperparams} reports the hyperparameters used for both models. Both are configured
conservatively given the relatively short and autocorrelated financial sample used in this thesis.
Restricting tree depth and, for Random Forest, requiring a minimum number of observations per leaf
prevent either model from fitting rules to a small number of unusual days rather than broader market
patterns. XGBoost builds trees sequentially, with each new tree correcting the errors of the
previous one; a low learning rate (0.05) limits how aggressively the model adapts to any single
tree's errors, while subsampling 80\% of rows and features per tree reduces overreliance on any
particular period or feature. These choices reflect a deliberately conservative configuration, given
the relatively short and autocorrelated financial sample used in this thesis and the general risk of
overfitting when flexible models are applied to limited historical data. A fixed random seed (42)
ensures the reported results are exactly reproducible.

\begin{table}[H]
\centering
\begin{tabular}{lll}
\toprule
\textbf{Hyperparameter} & \textbf{Random Forest} & \textbf{XGBoost} \\
\midrule
Number of trees        & 200 & 200  \\
Maximum depth          & 6   & 4    \\
Min. samples per leaf  & 5   & --   \\
Learning rate          & --  & 0.05 \\
Row subsample          & --  & 0.8  \\
Column subsample       & --  & 0.8  \\
Random seed            & 42  & 42   \\
\bottomrule
\end{tabular}
\caption{Hyperparameters used for the machine learning models.}
\label{tab:ml_hyperparams}
\end{table}

Figure~\ref{fig:xgb_tree} shows a single tree from an illustrative XGBoost fit, using the same
hyperparameters and the six-feature set above, trained once on the full training period for
visualisation purposes only and kept separate from the walk-forward evaluation used everywhere else
in this thesis.

\begin{figure}[H]
\centering
\includegraphics[width=\textwidth]{xgboost_tree_example.pdf}
\caption{A single tree (index 0 of 200) from an illustrative XGBoost fit on the training period.
Leaf values are that tree's additive contribution to the overall forecast; the final prediction sums
the contributions of all 200 trees, each individually small given the 0.05 learning rate.}
\label{fig:xgb_tree}
\end{figure}

Including the GARCH or EWMA forecasts as features at this stage would make the comparison circular:
any advantage attributed to machine learning could simply reflect access to the econometric model's
own output rather than a genuinely different source of predictive information. GARCH- and
EWMA-informed features are therefore introduced later, and only within the hybrid model
(Section~\ref{sec:hybrid}), where the machine learning component is explicitly tasked with modelling
the residual information left unexplained by GARCH rather than competing against it on equal terms.

All features are constructed exclusively from information available at the forecasting date, and
particular attention is paid to preventing data leakage during feature construction and model
training. The same features and forecasting target are used for both Random Forest and XGBoost to
ensure a homogeneous comparison.

Consistent with the walk-forward procedure described in Section~\ref{sec:oos}, the econometric
models are re-estimated daily, while Random Forest and XGBoost are re-trained weekly. Daily
re-training is inexpensive for the econometric models but unnecessary for the machine learning
models, whose more expensive fitting procedure does not require daily updating to remain reflective
of current market conditions.

\subsection{Volatility Forecast Evaluation}
% ============================================================
The volatility forecasts produced by all models are evaluated against the same five-day realised
volatility target. \\
The main evaluation metrics include Mean Absolute Error (MAE), Root Mean Squared Error (RMSE)
and the QLIKE loss function. MAE and RMSE are computed directly on the volatility scale:\\
\begin{equation}
\text{MAE} = \frac{1}{N}\sum_{t=1}^{N} \left| \hat{\sigma}_t - \sigma_t \right|,
\qquad
\text{RMSE} = \sqrt{\frac{1}{N}\sum_{t=1}^{N} \left( \hat{\sigma}_t - \sigma_t \right)^2},
\end{equation}
where $\sigma_t$ is the realised volatility target and $\hat{\sigma}_t$ the model forecast.
QLIKE, by contrast, is defined on the variance scale and penalises underprediction of risk more
heavily than overprediction of the same magnitude, which makes it particularly informative in a
risk-management context:
\begin{equation}
\text{QLIKE} = \frac{1}{N}\sum_{t=1}^{N} \left[ \frac{\sigma_t^2}{\hat{\sigma}_t^2} - \ln\left(\frac{\sigma_t^2}{\hat{\sigma}_t^2}\right) - 1 \right].
\end{equation}

These three measures are complementary rather than interchangeable, and each answers a slightly
different question about forecast quality. MAE expresses the average forecasting error directly in
percentage points of annualised volatility: a MAE of $0.049$, for example, means that a model's
forecasts deviate from realised volatility by approximately 4.9 percentage points on an average day.
RMSE is expressed on the same scale but squares errors before averaging, so it is more heavily
influenced by occasional large misses than by many small ones; a model whose RMSE is noticeably
higher than its MAE is one that is usually close to the target but occasionally misses it by a wide
margin, which is the typical pattern for volatility forecasts around sudden market moves. QLIKE
cannot be read in percentage-point terms, since it is a ratio-based loss defined on variance rather
than volatility, but its economic interpretation is more directly relevant to risk management: it
penalises a forecast that understates risk far more severely than one that overstates it by the same
amount. This asymmetry matters because, in a risk-management context, telling a decision-maker that
volatility will remain low shortly before a spike is a materially more costly error than being overly
cautious ahead of a calm period -- a distinction that MAE and RMSE, which treat over- and
under-prediction symmetrically, do not capture.

Reporting all three measures avoids relying on a single loss function, since no individual metric is
uniformly preferred for evaluating volatility forecasts \citep{poon2003}. These measures provide
complementary information about the accuracy of the forecasts and allow the performance of the
econometric, machine learning and hybrid models to be compared under the same criteria.

\subsection{Hybrid Model}\label{sec:hybrid}

The hybrid model combines the strongest econometric model considered in this thesis, GJR-GARCH,
with a machine learning component trained specifically to correct its forecasting errors. The
purpose of this design is not to combine model predictions mechanically, but to give each component
a clearly defined role: GJR-GARCH provides the baseline econometric volatility forecast, and the
machine learning component is trained to predict the residual information left unexplained by that
baseline.

The hybrid forecast takes the form

\begin{equation}
\hat{\sigma}^{Hybrid}_{t+1:t+5}
=
\hat{\sigma}^{GJR}_{t+1:t+5}
+
\hat{e}^{ML}_{t+1:t+5},
\end{equation}

where $\hat{\sigma}^{GJR}_{t+1:t+5}$ is GJR-GARCH's own five-day annualised volatility forecast and
$\hat{e}^{ML}_{t+1:t+5}$ is a machine learning correction trained to predict GJR-GARCH's forecasting
error. GJR-GARCH is used as the baseline rather than the symmetric GARCH(1,1) model because it is
the strongest-performing econometric model in this thesis (Section~\ref{sec:ml_comparison}), making
it the more demanding and more informative benchmark for evaluating whether machine learning adds
value on top of an already well-specified model.

XGBoost is used as the residual-correction component rather than Random Forest. This choice is
motivated by the structural similarity between gradient boosting and residual correction: XGBoost
builds each successive tree specifically to correct the errors of the trees before it, which is
mechanistically analogous to the task assigned to it here -- predicting the errors of a prior model
-- whereas Random Forest has no equivalent structural relationship to residual fitting.

The residual model is trained on the same six features used in the Research Question 1 comparison
(Section~\ref{sec:ml_features}), together with GJR-GARCH's own current forecast as a seventh feature.
Unlike in the Research Question 1 comparison, where including econometric model output as a feature
would have made the comparison circular, using GJR-GARCH's forecast as an input here is precisely
the point of a hybrid model: it allows the machine learning component to learn how GJR-GARCH's own
errors relate to its own forecast level, rather than only to raw return information.

Training residuals are obtained from a single GJR-GARCH fit on the full training period (2011--2022),
rather than a full walk-forward re-estimation, since this component's econometric behaviour has
already been separately validated (Section~\ref{sec:oos}) and does not need to be re-estimated here.
The residual model is trained on these in-sample (feature, residual) pairs. During out-of-sample
evaluation, the hybrid model reuses GJR-GARCH's already-validated walk-forward forecasts for
2023--2026 rather than recomputing them, and the XGBoost residual component is re-trained weekly,
consistent with the machine learning refit cadence used throughout this thesis. The hybrid model is
evaluated using exactly the same out-of-sample procedure, evaluation metrics, and backtesting
procedures as the individual models.

Because these training residuals come from a single in-sample fit, they are likely to be somewhat
smaller and less variable than the residuals GJR-GARCH actually produces out-of-sample. This is not
a source of look-ahead bias, since every training-period date precedes the test window entirely, but
it is a modelling simplification whose implications for the hybrid's calibration are discussed
alongside the results in Section~\ref{sec:hybrid_results}.

\subsection{Market Risk Measures}\label{sec:market_risk}

The volatility forecasts are subsequently used to estimate market-risk measures.

Value at Risk is calculated for the selected confidence levels using a common distributional
assumption (parametric, Gaussian) across all models. Expected Shortfall is also calculated to
measure the expected loss conditional on exceeding the VaR threshold.

Using the same assumptions for all models ensures that differences in the resulting risk measures
are primarily attributable to differences in the volatility forecasts rather than to differences in
the risk-measurement methodology.

\subsection{Backtesting}\label{sec:backtesting}

The VaR forecasts are evaluated using standard backtesting procedures.

The Kupiec test is used to evaluate unconditional coverage, assessing whether the observed
frequency of VaR violations is consistent with the nominal confidence level.

The Christoffersen test is used to evaluate the independence and conditional coverage of VaR
violations.

Both tests are formulated as hypothesis tests, and it is the failure to reject the null hypothesis
that is desirable for a well-specified model. The Kupiec test's null hypothesis is that the observed
proportion of VaR violations equals the nominal violation rate implied by the confidence level (e.g.
exactly 5\% of test days for a 95\% VaR); a high likelihood-ratio (LR) statistic and correspondingly
low $p$-value indicate that the violation rate is significantly too high or too low, and thus that
the model is miscalibrated. The Christoffersen test's null hypothesis is instead that violations are
independently distributed over time -- that a violation on one day does not make a violation on the
following day more likely; rejecting this null indicates that violations cluster together in time,
which is a more serious failure than an incorrect overall rate, since it implies the model performs
poorly precisely during extended periods of market stress rather than failing unpredictably.
Throughout this thesis, a $p$-value below the conventional 5\% threshold is treated as evidence
against the corresponding null hypothesis.

Because the volatility target and VaR forecasts in this thesis are constructed over a five-day
horizon while forecasts are produced daily, consecutive VaR windows overlap and share up to four of
their five underlying daily returns. This induces mechanical serial correlation in the violation
series that is unrelated to model quality: a single large market move can register as a violation
across several consecutive, overlapping forecast windows regardless of how well-specified the
underlying volatility model is. The Christoffersen independence test is therefore applied primarily
to a non-overlapping subsample of forecast dates, constructed so that no two consecutive observations
share any underlying daily return; the full-sample result, which reflects this structural overlap
rather than genuine violation clustering, is reported only as a secondary diagnostic.

Expected Shortfall forecasts are subjected to a simple backtest comparing the average realised
shortfall on violation days against the corresponding Expected Shortfall forecast.

The purpose of the backtesting exercise is to determine whether differences in volatility forecasting
accuracy translate into economically relevant differences in market-risk forecasting.

\subsection{Implementation and Reproducibility}
% ============================================================
All models -- econometric, machine learning, and, subsequently, the hybrid specification -- share a
common interface within the walk-forward procedure, so that every model is evaluated through
identical code rather than model-specific logic that could introduce inconsistencies between
families of models. Random Forest and XGBoost are both trained with a fixed random seed. An early
diagnostic run revealed that parallelised training introduced small but non-negligible
floating-point differences across otherwise identical XGBoost runs; XGBoost is therefore trained
single-threaded to guarantee exactly reproducible results. Random Forest is deterministic under a
fixed seed regardless of thread count, and is run single-threaded only to suppress excessive
diagnostic warning output from the underlying library. Full implementation details are documented in
the accompanying code repository, consistent with keeping the main document readable independently
of the underlying code.

\section{Results}
% ============================================================
\subsection{Econometric Model Comparison}
% ============================================================
Table~\ref{tab:econ_comparison} reports out-of-sample forecasting performance for the three
econometric models over the test period (3 January 2023 to 7 July 2026, 874 valid observations
after removing the tail rows for which the forward-looking target is undefined).

\begin{table}[H]
\centering
\begin{tabular}{lccc}
\toprule
\textbf{Model} & \textbf{MAE} & \textbf{RMSE} & \textbf{QLIKE} \\
\midrule
EWMA        & 0.04925 & 0.07656 & 0.43662 \\
GARCH(1,1)  & 0.04763 & 0.07197 & 0.38670 \\
GJR-GARCH   & 0.04358 & 0.06485 & 0.34217 \\
\bottomrule
\end{tabular}
\caption{Out-of-sample forecasting performance of the econometric models, 2023--2026.}
\label{tab:econ_comparison}
\end{table}

GJR-GARCH outperforms both GARCH(1,1) and EWMA on all three metrics simultaneously, with
improvements ranging from approximately 8\% to 12\% over GARCH(1,1) and 12\% to 22\% over EWMA
depending on the metric. The advantage is largest under QLIKE, consistent with GJR-GARCH's
asymmetric term allowing it to react more quickly to the volatility increases that typically follow
negative returns -- the well-documented leverage effect in equity index returns -- a pattern that
neither plain GARCH nor EWMA can capture. EWMA, despite having no fitted parameters and no
mean-reversion mechanism, remains within approximately 12\% of GJR-GARCH on MAE, confirming its
role as a reasonable naive benchmark rather than a negligible one.

These results should be interpreted as a first, partial answer to Research Question 1: among the
econometric models considered, explicitly modelling volatility asymmetry improves out-of-sample
forecasting accuracy. Whether this advantage persists, narrows, or is overturned once the machine
learning models are introduced -- and whether it translates into materially better market-risk
measures -- is assessed in the following subsections.

\subsection{Machine Learning Model Comparison}\label{sec:ml_comparison}
% ============================================================
Table~\ref{tab:full_comparison} extends the comparison to include Random Forest and XGBoost,
trained on the six-feature set described in Section~\ref{sec:ml_features} and evaluated over the
same out-of-sample test period and criteria as the econometric models.

\begin{table}[H]
\centering
\begin{tabular}{lccc}
\toprule
\textbf{Model} & \textbf{MAE} & \textbf{RMSE} & \textbf{QLIKE} \\
\midrule
EWMA          & 0.04925 & 0.07656 & 0.43662 \\
GARCH(1,1)    & 0.04763 & 0.07197 & 0.38670 \\
GJR-GARCH     & 0.04358 & 0.06485 & 0.34217 \\
Random Forest & 0.04341 & 0.07137 & 0.41497 \\
XGBoost       & 0.04457 & 0.07471 & 0.42450 \\
\bottomrule
\end{tabular}
\caption{Out-of-sample forecasting performance across all econometric and machine learning models, 2023--2026.}
\label{tab:full_comparison}
\end{table}

No single model dominates across all three criteria, though the margins differ substantially.
Random Forest achieves the lowest MAE of any model (0.04341), narrowly ahead of GJR-GARCH (0.04358)
by 0.39\% -- a gap small enough that its statistical significance is assessed formally in
Section~\ref{sec:dm} rather than treated as decisive here. XGBoost trails both, ranking third on
MAE. On RMSE and QLIKE, GJR-GARCH is unambiguously the best-performing model, ahead of the
next-best model (Random Forest on RMSE, GARCH(1,1) on QLIKE) by approximately 9\% and 12\%
respectively, and neither machine learning model improves on the symmetric GARCH(1,1) benchmark
under QLIKE.

These results provide a direct, though qualified, answer to the first part of Research Question 1.
Under the strictly out-of-sample, leakage-controlled framework adopted in this thesis, machine
learning offers at most a marginal advantage over the best econometric model on average
point-forecast accuracy, and no advantage once larger errors and tail risk are weighted more
heavily. This finding is consistent with \citet{poon2003}'s observation that forecast rankings are
sensitive to the evaluation criterion used, which extends naturally from the choice of
realised-volatility proxy they discuss to the choice between machine learning and econometric model
families. GJR-GARCH remains the strongest single model overall, ranking among the best two models
under every criterion considered. Whether the narrow Random Forest advantage on MAE is statistically
distinguishable from GJR-GARCH's performance is assessed via a Diebold-Mariano test in
Section~\ref{sec:dm}, once the hybrid model (Section~\ref{sec:hybrid}) has also been evaluated.

\subsection{Market Risk and Backtesting}
% ============================================================
Table~\ref{tab:kupiec_results} reports Kupiec unconditional coverage test results for all six
models at the 95\% and 99\% confidence levels. All models pass comfortably at 95\%. At 99\%, the
three econometric models remain comfortably calibrated ($p \geq 0.46$ in all three cases). Random
Forest is borderline, passing only marginally ($p=0.0535$, just above the conventional 5\%
threshold), while XGBoost and the hybrid model both fail outright -- the hybrid landing on exactly
XGBoost's violation count and statistics (17 violations; LR $=6.179$, $p=0.0129$), consistent with
the inherited-tail-weakness mechanism discussed in Section~\ref{sec:hybrid_results}.

\begin{table}[H]
\centering
\begin{tabular}{lccccc}
\toprule
Model & Confidence & Violation Rate (\%) & Nominal Rate (\%) & LR Statistic & $p$-value \\
\midrule
EWMA           & 95\% & 5.03 & 5.00 & 0.002 & 0.9629 \\
EWMA           & 99\% & 1.26 & 1.00 & 0.546 & 0.4601 \\
GARCH(1,1)     & 95\% & 4.23 & 5.00 & 1.138 & 0.2861 \\
GARCH(1,1)     & 99\% & 1.14 & 1.00 & 0.175 & 0.6754 \\
GJR-GARCH      & 95\% & 3.89 & 5.00 & 2.446 & 0.1178 \\
GJR-GARCH      & 99\% & 1.14 & 1.00 & 0.175 & 0.6754 \\
Random Forest  & 95\% & 5.84 & 5.00 & 1.221 & 0.2691 \\
Random Forest  & 99\% & 1.72 & 1.00 & 3.730 & 0.0535 \\
XGBoost        & 95\% & 6.06 & 5.00 & 1.957 & 0.1619 \\
XGBoost        & 99\% & 2.06 & 1.00 & 7.588 & 0.0059 \\
Hybrid         & 95\% & 5.38 & 5.00 & 0.256 & 0.6127 \\
Hybrid         & 99\% & 1.95 & 1.00 & 6.179 & 0.0129 \\
\bottomrule
\end{tabular}
\caption{Kupiec unconditional coverage test results across all models and confidence levels.}
\label{tab:kupiec_results}
\end{table}

Table~\ref{tab:es_backtest} compares the mean realised shortfall on violation days against the
corresponding Expected Shortfall forecast. Realised shortfall exceeds the forecast ES for every
model at both confidence levels, with the excess ranging from approximately 8\% to 24\% among the
standalone models, and reaching approximately 16--18\% for the hybrid model. This pattern is
consistent across all model families, which suggests it reflects the shared Gaussian distributional
assumption underlying the VaR and ES calculation (Section~\ref{sec:market_risk}) rather than a
weakness specific to any individual volatility forecast: financial returns are known to exhibit
heavier tails than the Gaussian distribution, so realised losses on the worst days tend to exceed
what a Gaussian-based Expected Shortfall anticipates. This is a known and expected limitation of the
modelling choice made in this thesis rather than an unexpected finding, and is discussed further in
the conclusions.

\begin{table}[H]
\centering
\begin{tabular}{lcccc}
\toprule
Model          & Confidence & Mean Realised Shortfall (\%) & Mean Forecast ES (\%) & Excess (\%) \\
\midrule
EWMA           & 95\% & 3.94 & 3.59 & 9.84 \\
EWMA           & 99\% & 6.05 & 4.98 & 21.46 \\
GARCH(1,1)     & 95\% & 4.11 & 3.81 & 7.89 \\
GARCH(1,1)     & 99\% & 6.44 & 5.50 & 17.02 \\
GJR-GARCH      & 95\% & 4.14 & 3.71 & 11.48 \\
GJR-GARCH      & 99\% & 6.40 & 5.65 & 13.34 \\
Random Forest  & 95\% & 3.88 & 3.29 & 17.97 \\
Random Forest  & 99\% & 5.36 & 4.33 & 23.90 \\
XGBoost        & 95\% & 3.83 & 3.29 & 16.61 \\
XGBoost        & 99\% & 5.22 & 4.26 & 22.45 \\
Hybrid         & 95\% & 3.88 & 3.34 & 16.21 \\
Hybrid         & 99\% & 5.24 & 4.43 & 18.11 \\
\bottomrule
\end{tabular}
\caption{Expected Shortfall backtest: mean realised shortfall vs. forecasted ES on violation days.}
\label{tab:es_backtest}
\end{table}

Figure~\ref{fig:forecast_error} shows the forecast error (actual minus forecast) for GJR-GARCH and
XGBoost over the test period. Both models record their largest single underprediction on the same
date, 2 April 2025, when realised annualised volatility reached 86.3\%: GJR-GARCH forecast 17.3\%
(an underprediction of 69.0 percentage points) and XGBoost forecast 13.3\% (an underprediction of
73.0 percentage points). The gap between the two models on this single worst day is modest, and
this episode alone is not strong evidence that GJR-GARCH copes better with extreme volatility than
XGBoost. The more reliable evidence for that claim is the aggregate pattern across the full test
period: XGBoost's QLIKE, Kupiec, and Expected Shortfall results are each independently weaker than
GJR-GARCH's, computed from all 874 out-of-sample observations rather than a single extreme event.

\begin{figure}[H]
\centering
\includegraphics[width=0.85\textwidth]{forecast_error_test.pdf}
\caption{Forecast error (actual minus forecast) for GJR-GARCH and XGBoost over the out-of-sample
test period. Positive values indicate underprediction of realised volatility. Both models record
their largest underprediction on 2 April 2025.}
\label{fig:forecast_error}
\end{figure}

As discussed in Section~\ref{sec:backtesting}, the full-sample Christoffersen independence test is
affected by mechanical serial correlation arising from the overlapping five-day forecast horizon and
is therefore not reported as the primary evidence of independence. Table~\ref{tab:christoffersen_nonoverlap}
reports the Christoffersen test applied instead to a non-overlapping subsample of forecast dates. All
six models pass comfortably ($p > 0.32$ in every case), indicating that, once the structural overlap
is removed, VaR violations do not cluster in time for any of the models considered.

\begin{table}[H]
\centering
\begin{tabular}{lcccc}
\toprule
Model          & $n$ & Violations & LR Statistic & $p$-value \\
\midrule
EWMA           & 175 & 9 & 0.982 & 0.3216 \\
GARCH(1,1)     & 175 & 7 & 0.587 & 0.4436 \\
GJR-GARCH      & 175 & 6 & 0.429 & 0.5126 \\
Random Forest  & 175 & 9 & 0.982 & 0.3216 \\
XGBoost        & 175 & 9 & 0.982 & 0.3216 \\
Hybrid         & 175 & 7 & 0.587 & 0.4436 \\
\bottomrule
\end{tabular}
\caption{Christoffersen independence test on a non-overlapping subsample (every 5th test date, 95\% confidence).}
\label{tab:christoffersen_nonoverlap}
\end{table}

\subsection{Hybrid Model Results}\label{sec:hybrid_results}
% ============================================================
Table~\ref{tab:hybrid_comparison} extends the model comparison to include the hybrid specification.

\begin{table}[H]
\centering
\begin{tabular}{lccc}
\toprule
\textbf{Model} & \textbf{MAE} & \textbf{RMSE} & \textbf{QLIKE} \\
\midrule
EWMA          & 0.04925 & 0.07656 & 0.43662 \\
GARCH(1,1)    & 0.04763 & 0.07197 & 0.38670 \\
GJR-GARCH     & 0.04358 & 0.06485 & 0.34217 \\
Random Forest & 0.04341 & 0.07137 & 0.41497 \\
XGBoost       & 0.04457 & 0.07471 & 0.42450 \\
Hybrid (GJR-GARCH + XGBoost) & 0.04019 & 0.06256 & 0.38484 \\
\bottomrule
\end{tabular}
\caption{Out-of-sample forecasting performance across all models, including the hybrid specification, 2023--2026.}
\label{tab:hybrid_comparison}
\end{table}

The hybrid model achieves the lowest MAE and RMSE of any model considered in this thesis, improving
on GJR-GARCH -- already the strongest standalone model -- by 7.8\% on MAE and 3.5\% on RMSE. On
QLIKE, however, it ranks second, behind GJR-GARCH and only marginally ahead of GARCH(1,1) (0.38484
versus 0.38670). This is the first indication that the hybrid's improvement in average point-forecast
accuracy does not extend uniformly to the loss function most directly relevant to risk management.

Table~\ref{tab:kupiec_results} confirms this directly. At 95\% confidence the hybrid passes the
Kupiec test comfortably. At 99\%, it fails -- by exactly the same margin as XGBoost (17 violations
out of 874; LR $=6.179$, $p=0.0129$) -- despite being built on GJR-GARCH, the best-calibrated model
in this thesis at this confidence level. Expected Shortfall backtesting shows the same
deterioration: excess over forecasted ES rises from 11.48\%/13.34\% for GJR-GARCH alone to
16.21\%/18.11\% for the hybrid, at 95\% and 99\% confidence respectively.

This outcome traces to a specific, identifiable mechanism. GJR-GARCH's forecasts run systematically
high over the test period (mean forecast 0.1425 against mean realised volatility 0.1285), and
the residual model correctly learns to shade this down: its correction is negative on 91.1\% of
days, with a mean adjustment of $-0.017$. This reduces symmetric error, which is why MAE and RMSE
improve. The same downward shift, however, increases the share of days on which the hybrid
underpredicts realised volatility from 32.7\% under GJR-GARCH alone to 46.3\% under the hybrid --
precisely the failure mode QLIKE penalises asymmetrically and that drives VaR violations. On the
days the 99\% VaR is actually breached, the correction is more negative still (mean $-0.022$) than
on an average day: the residual model shades hardest exactly where doing so is most costly.
Consistent with this mechanism, the hybrid's 17 violation dates at 99\% overlap with 15 of
XGBoost's and 16 of Random Forest's, but with only 9 of GJR-GARCH's own -- the hybrid's tail
behaviour resembles the machine learning models it inherited its correction from, not the
econometric model it was built on.

These results provide a complete, though qualified, answer to Research Question 2. Combining
GJR-GARCH with a machine learning residual correction improves average forecasting accuracy by a
larger margin than any standalone model achieves relative to any other. It does not improve, and in
this case actively worsens, calibration at the confidence level most relevant to market-risk
management, reproducing at the model-combination level the same pattern already observed for the
standalone machine learning models. This should not be read as a straightforward improvement: the
correct interpretation is that combining an econometric model with a machine learning correction can
produce the best available point forecasts while simultaneously producing worse risk-management
properties than the econometric model alone -- directly reinforcing this thesis's central finding
that improvements in standard forecast-accuracy metrics do not automatically translate into
improvements in market-risk measurement.

A methodological caveat is worth stating plainly. The residual model is trained on in-sample
residuals from a single GJR-GARCH fit over the training period, rather than on out-of-sample
walk-forward residuals, since GJR-GARCH's econometric behaviour was already separately validated
(Section~\ref{sec:oos}). In-sample residuals are typically smaller and less variable than genuine
out-of-sample errors, which would, if anything, be expected to bias the residual model toward
\emph{under}-correcting relative to what it would learn from harder, walk-forward-derived training
data. Since the hybrid's problem here is over-correction in the downward direction, this caveat does
not explain away the tail-calibration result; a residual model trained on genuinely out-of-sample
errors would plausibly shade forecasts down more aggressively still, not less.

\subsubsection{A Quantile-Loss Variant: Diagnosis and Limits}\label{sec:hybrid_quantile}

Motivated by the mechanism identified above -- a symmetric-loss residual correction improves
average accuracy at the cost of tail calibration -- a second hybrid variant was constructed to test
whether aligning the residual model's training objective with the actual risk-management target
would resolve this trade-off. Because Value at Risk at confidence level $\alpha$ is itself the
$(1-\alpha)$ quantile of the loss distribution, the residual model was retrained using a quantile
(pinball) loss, targeting an upper quantile of GJR-GARCH's in-sample residual distribution rather
than its conditional mean. All other aspects of the hybrid specification -- baseline model,
features, refit cadence -- are unchanged from Section~\ref{sec:hybrid}.

The quantile level was selected using only the training period (2011--2022), split chronologically
into an inner-training set (2011--2019) and an inner-validation set (2020--2022), with the test
period never consulted during selection. The selection criterion was revised once during this
process: an initial rule based solely on maximising the 99\% Kupiec $p$-value on the validation set
proved degenerate, since it is unbounded and would select an arbitrarily wide, uninformative
forecast purely because it is trivially well-calibrated. It was replaced with a rule requiring the
lowest validation-set MAE among candidates that also achieve acceptable 99\% coverage. This revision
was made using validation-period results only, before the test period was examined, but constitutes
a mild form of selection on validation data and is disclosed here for transparency. Under the
revised criterion, $\alpha=0.95$ was selected.

On the test period, this variant performs substantially worse than either GJR-GARCH alone or the
symmetric-loss hybrid, and is the worst-performing model considered anywhere in this thesis
(Table~\ref{tab:hybrid_quantile_comparison}). MAE more than doubles relative to GJR-GARCH (0.0972
versus 0.0436). The residual correction is now positive on every test day, and the resulting
forecast exceeds realised volatility on 94.16\% of test days (mean forecast 0.218 against mean
realised volatility 0.129). Its 99\% VaR fails calibration in the opposite direction to the
symmetric hybrid: 0.23\% observed violations against a 1\% nominal rate, rather than too many
(Table~\ref{tab:hybrid_quantile_kupiec}).

\begin{table}[H]
\centering
\begin{tabular}{lccc}
\toprule
Model & MAE & RMSE & QLIKE \\
\midrule
GJR-GARCH & 0.04358 & 0.06485 & 0.34217 \\
Hybrid (symmetric) & 0.04019 & 0.06256 & 0.38484 \\
Hybrid (quantile, $\alpha=0.95$) & 0.09715 & 0.11896 & 0.62505 \\
\bottomrule
\end{tabular}
\caption{Forecasting performance of the quantile-loss hybrid variant against GJR-GARCH and the symmetric-loss hybrid.}
\label{tab:hybrid_quantile_comparison}
\end{table}

\begin{table}[H]
\centering
\begin{tabular}{lccccc}
\toprule
Model & Confidence & Violation Rate (\%) & Nominal (\%) & LR & $p$-value \\
\midrule
GJR-GARCH           & 95\% & 3.89 & 5.00 & 2.446  & 0.1178 \\
GJR-GARCH           & 99\% & 1.14 & 1.00 & 0.175  & 0.6754 \\
Hybrid (symmetric)  & 95\% & 5.38 & 5.00 & 0.256  & 0.6127 \\
Hybrid (symmetric)  & 99\% & 1.95 & 1.00 & 6.179  & 0.0129 \\
Hybrid (quantile)   & 95\% & 0.80 & 5.00 & 49.359 & 0.0000 \\
Hybrid (quantile)   & 99\% & 0.23 & 1.00 & 7.633  & 0.0057 \\
\bottomrule
\end{tabular}
\caption{Kupiec test results for GJR-GARCH and both hybrid variants.}
\label{tab:hybrid_quantile_kupiec}
\end{table}

The cause is not an implementation defect but a distributional one. The inner-validation period
(2020--2022) used to select the quantile level has a substantially higher mean realised volatility
(0.1995) than the 2023--2026 test period (0.1285) -- a regime roughly 36\% calmer than the period
the correction was calibrated on. A 95th-percentile residual correction learned under COVID-era
market stress is simply too large for the comparatively calm conditions of the test period. This
illustrates a limitation more fundamental than a specific choice of hyperparameter: the residual
distribution this correction targets is not stationary across time, so a quantile level validated
on one historical regime does not necessarily transfer to another.

This result is reported as a negative finding rather than omitted, because it is informative in its
own right. Aligning a model's training objective with the true downstream use case is a necessary
but not sufficient condition for a well-calibrated hybrid model; it does not by itself protect
against non-stationarity in the quantity being modelled. Only the symmetric-loss hybrid improves on
GJR-GARCH's point-forecast accuracy, and it does so at a cost to tail-risk calibration; the quantile
variant improves on neither dimension, performing worse than GJR-GARCH on every forecasting metric
considered while also failing to achieve better-calibrated tail risk. Neither variant improves on
GJR-GARCH's own 99\% Kupiec performance, which remains the best-calibrated result in this thesis
($p=0.6754$) throughout. GJR-GARCH alone, without either machine learning extension, remains the
preferred model wherever calibrated tail-risk measurement is the primary objective, notwithstanding
the point-forecast accuracy gains available from the symmetric-loss hybrid.

\subsection{Statistical Significance of Forecasting Differences}\label{sec:dm}
% ============================================================
The comparisons in Sections~\ref{sec:ml_comparison} and \ref{sec:hybrid_results} identify
differences in MAE, RMSE and QLIKE across models, but do not by themselves establish whether these
differences are statistically distinguishable from zero. This is assessed using the
Diebold-Mariano test, applied to the two comparisons in this thesis where the ranking is close
enough to warrant formal testing: Random Forest against GJR-GARCH, and the symmetric-loss hybrid
against GJR-GARCH.

For a given loss function, let $d_t = L(\hat\sigma_t^{(a)}, \sigma_t) - L(\hat\sigma_t^{(b)}, \sigma_t)$
denote the loss differential between models $a$ and $b$ at date $t$. The Diebold-Mariano test
evaluates the null hypothesis of equal expected loss, $H_0: E[d_t] = 0$. Because forecasts in this
thesis are constructed over a five-day horizon while produced daily, consecutive loss differentials
overlap and share underlying return information -- the same structural feature that required a
non-overlapping subsample for the Christoffersen independence test (Section~\ref{sec:backtesting}).
Rather than subsampling, which would discard most of the test period for a test of this kind, the
statistic here uses a Newey-West heteroskedasticity- and autocorrelation-consistent (HAC) variance
estimator with a Bartlett kernel at four lags, matching the overlap length implied by the five-day
horizon. This choice affects the strength, though not the binary significance verdict, of two
comparisons: GJR-GARCH versus Random Forest and GJR-GARCH versus the hybrid model, both on RMSE,
sit just above the 5\% threshold under the naive i.i.d.\ assumption ($p=0.054$ in both cases) and
are pushed substantially further from significance once the overlap-induced autocorrelation is
properly accounted for ($p=0.215$ and $p=0.163$ respectively). The HAC correction therefore does not
reverse any conclusion reached under the naive assumption, but it does show that the apparent
borderline evidence in these two cases was considerably weaker than a naive analysis would suggest,
reinforcing rather than merely permitting the use of a robust variance estimator here.

Table~\ref{tab:dm_results} reports the results. Random Forest's narrow MAE advantage over GJR-GARCH
(Section~\ref{sec:ml_comparison}) is not statistically significant ($p=0.910$): the 0.4\% gap
between the two models is consistent with sampling noise rather than a genuine difference in
forecasting ability. GJR-GARCH's RMSE advantage over Random Forest is also not significant
($p=0.215$), but its QLIKE advantage is ($p=0.016$). For the hybrid comparison, the hybrid's MAE
improvement over GJR-GARCH is statistically significant ($p=0.005$) -- the one genuinely confirmed
forecasting gain identified anywhere in this thesis -- while its RMSE improvement is not
($p=0.163$). GJR-GARCH's QLIKE advantage over the hybrid is significant ($p=0.029$), confirming
that the hybrid's tail-calibration cost, documented mechanistically in Section~\ref{sec:hybrid_results},
is a statistically genuine effect rather than an artefact of a single unusual period.

\begin{table}[H]
\centering
\begin{tabular}{llcccl}
\toprule
Comparison & Loss & DM statistic & $p$-value & Lower loss & Significant (5\%) \\
\midrule
GJR-GARCH vs.\ Random Forest & MAE   & 0.113  & 0.910 & Random Forest & No \\
GJR-GARCH vs.\ Random Forest & RMSE  & $-$1.241 & 0.215 & GJR-GARCH     & No \\
GJR-GARCH vs.\ Random Forest & QLIKE & $-$2.417 & 0.016 & GJR-GARCH     & Yes \\
GJR-GARCH vs.\ Hybrid        & MAE   & 2.808  & 0.005 & Hybrid        & Yes \\
GJR-GARCH vs.\ Hybrid        & RMSE  & 1.396  & 0.163 & Hybrid        & No \\
GJR-GARCH vs.\ Hybrid        & QLIKE & $-$2.184 & 0.029 & GJR-GARCH     & Yes \\
\bottomrule
\end{tabular}
\caption{Diebold-Mariano test results (Newey-West HAC variance, Bartlett kernel, 4 lags). $n=874$ for all comparisons.}
\label{tab:dm_results}
\end{table}

These results sharpen the answers to both research questions. On Research Question 1, the apparent
machine learning advantage identified in Section~\ref{sec:ml_comparison} does not survive formal
testing: Random Forest's MAE ranking is not statistically distinguishable from GJR-GARCH's, while
GJR-GARCH's QLIKE advantage is. On Research Question 2, the hybrid model's improvement in average
forecasting accuracy is the one confirmed statistically significant gain identified anywhere in this
thesis, but it comes at a significant, not merely apparent, cost in tail-risk-relevant accuracy.

\section{Conclusions}
% ============================================================

This thesis set out to evaluate whether machine learning methods improve the out-of-sample
forecasting of financial market volatility relative to traditional econometric models, whether any
such improvement translates into more accurate market-risk measures, and whether combining
econometric and machine learning approaches offers additional value over either alone.

On forecasting accuracy, the initial evidence for a machine learning advantage does not survive
formal testing. Random Forest's apparent MAE edge over GJR-GARCH is not statistically significant
(Diebold-Mariano $p=0.910$; Section~\ref{sec:dm}): despite producing the lowest MAE among the five
standalone models, this ranking is consistent with sampling noise rather than a genuine forecasting
advantage. GJR-GARCH's advantage on QLIKE, by contrast, is statistically significant against both
machine learning models. The honest answer to the first part of Research Question 1 is therefore
that machine learning offers no statistically demonstrable improvement in out-of-sample volatility
forecasting over a well-specified econometric benchmark in this setting, and shows a significant
disadvantage under the criterion most relevant to risk management.

The market-risk results reinforce this conclusion. At 99\% confidence, XGBoost fails the Kupiec test
outright, and Random Forest passes only marginally ($p=0.0535$), while all three econometric models
remain comfortably calibrated ($p \geq 0.46$). That both machine learning models show weaker tail
calibration than every econometric model at this confidence level -- one failing outright, the
other only narrowly avoiding the same fate -- suggests a systematic weakness in tail-risk
calibration associated with the tree-based approach considered here, consistent with these models'
inability to extrapolate parametrically beyond previously observed volatility levels. Once the
mechanical serial correlation induced by the overlapping five-day forecast horizon is accounted for
via a non-overlapping subsample, none of the five standalone models exhibits genuine clustering of
VaR violations. Where all five do share a common weakness is in the magnitude of losses on days VaR
is breached: realised shortfall exceeds forecasted Expected Shortfall for every model at both
confidence levels, with excess ranging from approximately 8\% to 24\%. Because this pattern is
uniform across model families, it is best attributed to the shared Gaussian distributional
assumption used throughout this thesis rather than to any individual forecast; a fat-tailed
alternative, such as a Student-$t$ distribution, would be expected to narrow, though not necessarily
eliminate, this gap.

Research Question 2 asked whether combining GJR-GARCH with a machine learning residual correction
adds value over either component alone. The answer is genuinely two-sided, and this is the central
finding of this thesis. The hybrid model achieves the lowest MAE and RMSE of any model considered,
and its MAE improvement over GJR-GARCH is the one statistically significant forecasting gain
identified anywhere in this study (Diebold-Mariano $p=0.005$). At the same time, it fails the 99\%
Kupiec test by exactly the same margin as XGBoost, and this deterioration traces to a specific,
demonstrated mechanism: the residual correction learns a near-uniform downward shift, appropriate on
average but applied indiscriminately, that shades hardest exactly on the days it is most costly to
do so. The hybrid's tail-risk violations overlap overwhelmingly with the machine learning models'
own violations rather than with GJR-GARCH's, and GJR-GARCH's QLIKE advantage over the hybrid is
itself statistically significant. A subsequent attempt to correct this weakness directly, by training
the residual model with a quantile loss aligned with the definition of Value at Risk itself, made the
problem worse rather than better: the resulting model was the least accurate considered anywhere in
this thesis, its calibration failure traced to non-stationarity in the residual distribution between
the validation period used to select it and the test period.

Taken together, this thesis's central finding is that improvements in standard point-forecast
accuracy do not translate automatically into improvements in market-risk measurement, and in the
case of the hybrid model considered here, a statistically confirmed forecasting improvement is
accompanied by a statistically confirmed deterioration in the one property -- tail-risk calibration
-- that this thesis's second half exists to evaluate. GJR-GARCH alone, without either machine
learning extension, remains the best-calibrated model in this thesis at the confidence level most
relevant to market risk management, and no model or model combination considered here improves on
it in that specific respect.

This analysis is subject to several limitations. It is based on a single asset over a single
out-of-sample window, and the extent to which these results generalise to other assets, markets, or
periods is not established here. The market-risk measures rely on a parametric Gaussian
distributional assumption, which is shown to systematically understate realised shortfall; a
fat-tailed alternative was considered outside the scope of this thesis but represents a natural
extension. The residual model underlying both hybrid variants is trained on in-sample residuals from
a single GJR-GARCH fit rather than on genuinely out-of-sample walk-forward residuals; this is not a
source of look-ahead bias, since every training-period date precedes the test window entirely, but
it is a simplification whose robustness to a harder, walk-forward-derived training signal was not
tested. Finally, the quantile-loss hybrid's failure illustrates a limitation more general than any
single hyperparameter choice: aligning a model's training objective with its true downstream use is
necessary but not sufficient when the quantity being modelled is not stable across time, and any
future extension in this direction would need to address that non-stationarity directly rather than
through a fixed, validation-period-selected correction.

\subsection{Directions for Future Research}
% ============================================================
Several extensions follow directly from findings in this thesis, rather than from a general
wish to increase complexity.

\begin{itemize}
    \item \textbf{Regime-aware residual correction.} The quantile-loss hybrid variant failed
    because a correction calibrated on one volatility regime did not transfer to a calmer one
    (Section~\ref{sec:hybrid_quantile}). A natural extension is a residual correction that
    conditions explicitly on the current volatility regime -- for example, using recent realised
    volatility relative to its longer-run average as an additional input -- rather than applying
    a single fixed correction learned once.

    \item \textbf{Fat-tailed distributional assumptions.} The Gaussian assumption used throughout
    this thesis systematically understates realised shortfall (Section~\ref{sec:market_risk}). A
    Student-$t$ or other fat-tailed distribution for VaR and Expected Shortfall would test whether
    this gap narrows, and by how much, without requiring any change to the underlying volatility
    forecasts.

    \item \textbf{Alternative hybrid architectures.} This thesis combines GJR-GARCH and XGBoost
    additively, with the machine learning component predicting GJR-GARCH's residual error. A
    weighted or learned combination of the two models' independent forecasts, rather than a
    residual correction applied to one baseline, represents a structurally different approach
    worth comparing directly against the results reported here.

    \item \textbf{Walk-forward residual training.} The hybrid's residual model is trained on a
    single in-sample fit rather than genuinely out-of-sample residuals (Section~\ref{sec:hybrid}).
    Repeating the analysis with walk-forward-derived training residuals would test the robustness
    of the tail-calibration finding to this simplification.

    \item \textbf{Generalisation beyond a single asset and period.} All results in this thesis are
    specific to the S\&P 500 over 2023--2026. Replicating the comparison on other equity indices,
    other asset classes, or other out-of-sample windows would establish whether the disconnect
    between forecasting accuracy and risk calibration identified here is a general property of
    this modelling approach or specific to the conditions studied.
\end{itemize}
\section*{Declaration of AI Use}
% ============================================================
This thesis was developed with assistance from Claude (Anthropic), used in two
capacities: as a conversational assistant for methodological discussion, drafting and revising
this document, and reviewing code and results for correctness; and, via Claude Code, for
implementing the data pipeline, econometric and machine learning models, and statistical
analyses described, under the author's direct instruction and review at every stage. This
iterative process included several rounds of auditing that identified and corrected genuine
implementation errors -- including feature-timing inconsistencies, forecast staleness in the
walk-forward evaluation, and inaccuracies in reported statistics -- before they affected the
results presented here. All code and results created using AI were reviewed and verified by the author.
\bibliographystyle{apalike}
\bibliography{references}

\end{document}